%% =====================================================================
%%  第5章 数値実験
%%  書くこと：実験の目的，設定（再現に必要な情報），結果，考察
%%  以下の実験内容と数値は書き方の例（数値は架空）．自分の実験に合わせて書き換える．
%% =====================================================================
\chapter{数値実験}
\label{chap:experiments}

\begin{writingnote}
以下は記述・組版の例であり，実験設定と数値は説明用に作成した架空のものである．
図表の平均・標準偏差も実測値ではない．自分の実験内容と結果に置き換えること．
\end{writingnote}

本章では，提案手法の有効性を検証するために行った数値実験について述べる．
\begin{writingnote}
最初に，実験で何を確かめるのか（実験の目的や問い）を，「はじめに」の章で述べた目的と対応させて書く．
確かめることが複数あるときは，それぞれを番号付きの問いや実験として示し，結果と考察も同じ順に述べるとよい．
\end{writingnote}

\section{実験設定}
\label{sec:setup}

\begin{writingnote}
他の人が同じ実験を再現できるように，問題例（データ），比較手法，評価指標，パラメータ，実行環境を明記する．
問題例と比較手法は，選んだ理由も書く（例：既存研究で広く使われている問題例である，代表的な既存手法である）．
比較が公平になるように実験を計画する．確かめる点は，ソースファイルのこの位置にコメントとして書いてある．
\end{writingnote}

%% 実験を計画するときに確かめる点
%% ・調整用・検証用のデータと，最終評価用のデータを分ける．最終評価の結果で調整しない．
%% ・各手法には同じテスト問題例を用いる．計算時間・反復数などの予算と停止条件を明記する．
%% ・乱数を使う手法は複数の種で実行し，種の一覧，問題例数，問題例ごとの実行回数を書く．
%% ・問題例間のばらつきと，同じ問題例に対する乱数の種によるばらつきを区別する．
%%   平均・標準偏差がどちらを集計した値かを，表や図の説明にも明記する．
%% ・調整するパラメータの探索範囲・試行数・選択基準を記録する．

%% \paragraph{...} は目次に載らない小見出し．
\paragraph{問題例}
品物数 $n \in \{100, 200, \dots, 1000\}$ の 0-1 ナップサック問題の問題例を，$n$ ごとに 10 個ずつ生成した．
品物の価値 $v_i$ と重さ $w_i$ は 1 以上 100 以下の整数の一様乱数とし，容量 $W$ は重さの合計の半分とした．

\paragraph{比較手法}
動的計画法による厳密解法~\cite{kellerer2004knapsack}と，\cref{sec:greedy}で述べた貪欲法を比較手法とした．

\paragraph{評価指標}
最適値との差（ギャップ）と計算時間を評価指標とした．
ギャップは，最適値を $z^{*}$，得られた解の目的関数値を $z$ として，$100 (z^{*} - z) / z^{*}$ [\%] で計算した．

\paragraph{パラメータと実行環境}
%% 表は table 環境で書き，キャプションは表の「上」に付ける．
%% 罫線は \toprule・\midrule・\bottomrule（booktabs）を使い，縦罫線は使わない．
提案手法のパラメータを\cref{tab:parameters}に示す．
ペナルティの重み $\lambda$ は Optuna~\cite{akiba2019optuna} を用いて調整した．
実装には Python~\cite{python} を用い，CPU が○○，メモリが○○~GB の計算機で実験を行った\footnote{%
  実行環境の記載例：OS ○○，Python ○○，ライブラリ○○（version ○○）．}．

\begin{table}[tbp]
  \centering
  \caption{提案手法のパラメータ（架空の設定例）}
  \label{tab:parameters}
  \begin{tabular}{lc}
    \toprule
    パラメータ & 値 \\
    \midrule
    ペナルティの重み $\lambda$ & 10 \\
    解の改善の反復回数 & 1000 \\
    乱数の種 & 0 \\
    \bottomrule
  \end{tabular}
\end{table}

\section{実験結果}
\label{sec:results}

\begin{writingnote}
図表を示すだけでなく，本文で図表を参照し，そこから読み取れる全体の傾向と主な結果を文章で述べる．図表の数値をすべて文章で繰り返す必要はない．
予想と異なる結果や，提案手法が劣る結果も省かずに示す．
手法間の差を述べるときは，その差がばらつき（標準偏差など）に比べて十分に大きいかを確かめる．
結果の解釈（なぜその結果になったのか）は，次の「考察」の節で述べる．
\end{writingnote}

%% 図は figure 環境で書き，キャプションは図の「下」に付ける．図のファイルは figures/ に置く．
%% グラフは PDF（ベクター形式）で保存すると，拡大しても粗くならない．
%% 図を並べるときは subfigure 環境を使う．\cref で「図5.1(a)」のように参照できる．
各手法のギャップと計算時間を\cref{fig:results}に示す．
\cref{fig:results-gap}より，提案手法のギャップはすべての $n$ で貪欲法より小さい．
また，\cref{fig:results-time}より，厳密解法の計算時間は $n$ とともに急激に増加するのに対し，
提案手法の計算時間はほぼ $n$ に比例して緩やかに増加する．

\begin{figure}[htbp]
  \centering
  \begin{subfigure}{0.48\linewidth}
    \centering
    \includegraphics[width=\linewidth]{results-gap.pdf}
    \caption{最適値とのギャップ}
    \label{fig:results-gap}
  \end{subfigure}
  \hfill
  \begin{subfigure}{0.48\linewidth}
    \centering
    \includegraphics[width=\linewidth]{results-time.pdf}
    \caption{計算時間}
    \label{fig:results-time}
  \end{subfigure}
  \caption[品物数と各手法の性能（架空データ）]{品物数 $n$ と各手法の性能（架空データ：各 $n$ で 10 個の問題例の平均を想定）}
  \label{fig:results}
\end{figure}

$n = 1000$ のときの結果を\cref{tab:results}に示す．各列の最良値を太字で示している．

\begin{table}[htbp]
  \centering
  \caption[$n = 1000$ における結果（架空データ）]{$n = 1000$ における各手法の結果（架空データ：10 個の問題例の平均 $\pm$ 標準偏差を想定）}
  \label{tab:results}
  \begin{tabular}{lrr}
    \toprule
    手法 & ギャップ [\%] & 計算時間 [s] \\
    \midrule
    厳密解法 & $\mathbf{0.00 \pm 0.00}$ & $83.0 \pm 21.7$ \\
    貪欲法   & $1.07 \pm 0.35$ & $\mathbf{0.01 \pm 0.00}$ \\
    提案手法 & $0.22 \pm 0.08$ & $2.00 \pm 0.12$ \\
    \bottomrule
  \end{tabular}
\end{table}

\section{考察}
\label{sec:discussion}

\begin{writingnote}
結果の解釈を述べる．結果を繰り返すだけにならないように注意する．
主な結果ごとに，なぜその結果になったのかを説明する．推測にとどまる説明は「〜と考えられる」のように書き，実験で確かめた事実と区別する．
既存研究で報告されている結果と比べられる場合は，一致する点と異なる点を述べる．
そのうえで，どのような条件で提案手法が有効なのか，どのような場合には有効でないのか（限界）を，結果を根拠に議論する．
ここで述べた限界は，「結論」の章の今後の課題につながる．
\end{writingnote}

本実験の条件では，……について……という傾向が見られた．
一方，……では……であり，適用には……という制約がある．
